\chapter{Trigonometria e Funções Trigonométricas}

Uma importante parte da Matemática é o estudo dos ângulos e das relações entre ângulos e distâncias. Para ilustrar essa importância utilizaremos o artigo a seguir que exemplifica como os conceitos da Trigonometria podem ser utilizados para resolver problemas da vida real.


\subsection{Eratóstenes e a circunferência da Terra}

\begin{flushright}
JOSÉ ROBERTO V. COSTA
\end{flushright}

Eratóstenes nasceu em Cirene, Grécia, e morreu em Alexandria, Egito, no terceiro século AEC. Ele era bibliotecário-chefe da famosa Biblioteca de Alexandria, e foi lá que ele encontrou, num velho papiro, indicações de que ao meio-dia de cada 21 de junho na cidade de Assuã (ou Syene, no grego antigo) 800 km ao sul de Alexandria, uma vareta fincada verticalmente no solo não produzia sombra.

Cultura inútil, diriam alguns. Não para um homem observador como Eratóstenes. Ele percebeu que o fenômeno não ocorria no mesmo dia e horário em Alexandria, e pensou:

— Se o mundo é plano como uma mesa, então as sombras das varetas têm de ser iguais. E se isto não acontece é porque a Terra deve ser curva!

\begin{figure}[h]
	\begin{center}
		\includegraphics[width=0.8\textwidth,clip]{siena}
	\end{center}
\end{figure}

Mais do que isso. Quanto mais curva fosse a superfície da Terra, maior seria a diferença no comprimento das sombras.

O Sol deveria estar tão longe que seus raios de luz chegam à Terra paralelos. Varetas fincadas verticalmente no chão em lugares diferentes lançariam sombras de comprimentos distintos.

Eratóstenes decidiu fazer um experimento. Ele mediu o comprimento da sombra em Alexandria ao meio-dia de 21 de junho, quando a vareta em Assuã, ao Sul do Egito, não produzia sombra. Assim, ele obteve o ângulo A, conforme a figura abaixo.

\begin{figure}[H]
	\begin{center}
		\includegraphics[height=4cm,width=0.4\textwidth,clip]{siena2}
		 \quad \quad
		\includegraphics[height=4cm,width=0.2\textwidth,clip]{siena3}
	\end{center}
\end{figure}

Preste atenção na figura acima. O ângulo \textbf{B} terá o mesmo valor de \textbf{A}, pois o desenho do experimento de Eratóstenes se reduz a uma geometria muito simples: \textit{se duas retas paralelas interceptam uma reta transversal, então os ângulos correspondentes são iguais. }

As retas paralelas são os raios de luz do Sol e a reta transversal é a que passa pelo centro da Terra e pela vareta em Alexandria. O ângulo B (também igual a $7^{\circ}$), é a uma fração conhecida da circunferência da Terra e corresponde à distância entre Assuã e Alexandria!

Eratóstenes sabia que essa distância valia cerca de 800 km e então pensou: $7^{\circ}$ são aproximadamente 1/50 de uma circunferência ($360^{\circ}$). E isso corresponde a cerca de 800 km. Oitocentos quilômetros vezes cinquenta são \textbf{quarenta mil quilômetros}, de modo que deve ser este o valor da circunferência da Terra.

VALOR ENCONTRADO ATUALMENTE: cerca de 40.072 km ao longo da linha do equador. Um erro muito pequeno para uma medida tão simples, e feita há tanto tempo! Com a circunferência, podemos calcular o diâmetro e o raio ou ainda o volume e a área da superfície, através de fórmulas simples.

Repare que o conhecimento utilizado por Eratóstenes (retas paralelas cortadas por uma transversal) é formalmente adquirido hoje nas aulas de geometria do Ensino Fundamental. \\




Fonte: Costa, J. R. V. Eratóstenes e a circunferência da Terra. Astronomia no Zênite, jul 2000. Disponível em: http://www.zenite.nu/eratostenes-e-a-circunferencia-da-terra/

\section{O conceito de Ângulo}



Um \textbf{ângulo} é determinado por uma abertura entre duas retas de mesma origem. Para medirmos ângulos precisamos estabelecer o valor do ângulo total (uma volta completa no sentido anti-horário) e os demais ângulos são medidos de maneira proporcional. Existem duas formas usuais de medirmos ângulos: definindo o ângulo total como $360$ \textbf{graus} ou como $2\pi$ \textbf{radianos}. Vejamos algumas medições em cada um dos dois sistemas:

\begin{figure}[h]
	\begin{center}
	\includegraphics[width=0.3\textwidth,trim={0 0 10,5cm 1cm},clip]{circ1}
		\hspace{1cm}
	\includegraphics[width=0.3\textwidth,trim={0 0 10,5cm 1cm},clip]{circ2}
	\end{center}
\end{figure}

Considere um ângulo $\theta$. Denotamos por $\theta_g$ a medida desse ângulo em graus e $\theta_r$ em radianos. Esses dois valores estão relacionados pela equação\footnote{Isso quer dizer que podemos obter $\theta_g$ se conhecemos $\theta_r$ ou vice versa.}
\[\frac{360}{2\pi}=\frac{\theta_g}{\theta_r}.\]

A unidade de medida padrão no ensino superior é radianos e para transformar para radianos é comum usar a relação equivalente $180^{\circ}=\pi rad$, conforme ilustra o exemplo que segue.

\begin{exer}
Converta $120^{\circ}$ para radianos. \\
\color{red}
Para fazermos essa conversão usamos o fato de que $180^{\circ}$ é equivalente a $\pi rad$ e montamos a seguinte regra de três.
$$\begin{array}{cccccc}
\text{graus} &&& \text{radianos}   \\
180^{\circ} &&& \pi rad  \\
 120^{\circ} &&& x
\end{array}$$
Assim, $180x=120\pi \Longrightarrow x=\frac{2 \pi}{3}rad$. Logo, $120^{\circ}$ é equivalente a $\frac{2 \pi}{3}rad$.
\color{black}
\end{exer}
Pode ser que seja necessário voltar de radianos para graus. O exemplo seguinte ilustra essa situação.
\begin{exe}
Converta $5\pi/3$ radianos para graus. \\

\textbf{Solução:}
\color{red}
Para converter de radianos para graus, basta usarmos a relação $180^{\circ}= \pi rad$. Assim,

$$5\pi/3 rad = 5 \times \frac{180^{\circ}}{3}= 300^{\circ}.$$
\end{exe}

 \section{Relações no Triângulo Retângulo}

 Importantes relações são estabelecidas entre os ângulos e os lados de triângulos retângulos. Considere um triângulo retângulo que contém um ângulo $0<\theta<\pi/2$ conforme figura abaixo:

    \begin{figure}[h]
  	\begin{center}
  		\includegraphics[width=0.5 \textwidth,trim={1cm 20cm 19cm 9cm} ,clip]{figdeflados}
  	\end{center}
  \end{figure}

 Em relação ao ângulo $\theta$, os lados desse triângulo recebem os seguintes nomes:
 \begin{itemize}
 	\item cateto oposto
 	\item cateto adjacente
 	\item hipotenusa
 \end{itemize}

 Como a soma dos ângulos internos de qualquer triângulo é $\pi$ radianos, o terceiro ângulo (que chamaremos de $\alpha$) desse triângulo mede $\pi/2-\theta$ (pois $\alpha+\theta+\pi/2=\pi$). Portanto, o cateto oposto ao ângulo $\theta$ é o cateto adjacente ao ângulo $\alpha$ e o cateto adjacente ao ângulo $\theta$ é o cateto oposto ao ângulo $\alpha$.

 Estudaremos as relações entre os lados de dois triângulos retângulos que possuem um mesmo ângulo $\theta$. Fixamos dois desses triângulos e nomeamos os lados desses triângulos da seguinte forma:

  \begin{figure}[h]
 	\begin{center}
 		\includegraphics[width=0.5 \textwidth,trim={1cm 28cm 17cm 1cm} ,clip]{figdeffunc}
 	\end{center}
 \end{figure}

Como ambos os triângulos possuem um ângulo reto e um ângulo $\theta$, o terceiro ângulo desses dois triângulos mede $\pi/2-\theta$. Ou seja, eles são semelhantes pelo caso $AAA$ (ângulo ângulo ângulo) e podemos estabelecer as seguintes relações:

\[\frac{y_1}{z_1}=\frac{y_2}{z_2}, \quad \frac{x_1}{z_1}=\frac{x_2}{z_2}, \quad \frac{y_1}{x_1}=\frac{y_2}{x_2}.\]

Ou seja, independente dos valores dos catetos e da hipotenusa, as razões entre os lados de um triângulo retângulo que contém um ângulos $\theta$ são constantes. Essas constantes (que dependem de $\theta$) recebem nomes, os principais deles são:\footnote{Outras razões recebem os nomes de secante, cosseno e cotangente.}
\begin{description}
	\item \textbf{Seno de theta:} $\ds\sen\theta=\frac{\mbox{cateto oposto}}{\mbox{hipotenusa}}$.
	\item \textbf{Cosseno de theta:} $\ds\cos\theta=\frac{\mbox{cateto adjacente}}{\mbox{hipotenusa}}$.
	\item \textbf{Tangente de theta:} $\ds\tg\theta=\frac{\mbox{cateto oposto}}{\mbox{cateto adjacente}}$.
\end{description}

Assim definidos, seno, cosseno e tangente são funções cujo domínio é o intervalo $(0,\pi/2)$. Para estendermos essas definições e obtermos funções com domínio real precisaremos generalizar a noção de ângulos para valores maiores que $2\pi$ e ângulos negativos e posteriormente, estender a noção de seno, cosseno e tangente desses ângulos através do círculo trigonométrico unitário.

\begin{obs} É possível calcular os valores de seno, cosseno e tangente para os ângulos $\pi/4$, $\pi/3$ e $\pi/6$ utilizando a relação de Pitágoras no quadrado de lado $1$ e no triângulo equilátero de lado $1$. Esses ângulos são chamados de notáveis porque para demais ângulos não existe uma maneira tão simples de estabelecer os valores de seno, cosseno e tangente.
\end{obs}

\begin{exer}
	Calcule os valores de seno, cosseno e tangente para os ângulos $\pi/4$, $\pi/3$ e $\pi/6$ utilizando a relação de Pitágoras no quadrado de lado $1$ e no triângulo equilátero de lado $1$, justificando a tabela a seguir. \\
	\color{red}
	Vamos mostrar as razões trigonométricas para o ângulo de $45^{\circ}$. Para isso, considere um quadrado de lado 1 e tracemos sua diagonal. Ao fazermos esse procedimento, dividimos o quadrado em dois triângulos retângulos iguais como mostra a fogura \ref{quadrado}
	\begin{figure}[h!]
	    \centering
	    \includegraphics[scale=0.4]{quadrado}

	    \caption{Quadrado de lado 1}
	    \label{quadrado}
	\end{figure}
	Podemos observar que :
	\begin{itemize}
	    \item Cateto oposto a $45^{\circ}$ vale 1;
	    \item Cateto adjascente a $45^{\circ}$ vale 1;
	    \item Hipotenusa desses triângulos retângulos é $\sqrt{2}$.
	\end{itemize}
	Assim,
	\begin{itemize}
	    \item $\sin(45^{\circ})= \frac{1}{\sqrt{2}}= \frac{\sqrt{2}}{2}$;
	     \item $\cos(45^{\circ})= \frac{1}{\sqrt{2}}= \frac{\sqrt{2}}{2}$;
	      \item $\tan(45^{\circ})= \frac{1}{1}= 1$.
	\end{itemize}
	Para encontrarmos as razões trigonométricas para os ângulos de $30^{\circ}$ e $60^{\circ}$ basta traçarmos uma altura (que é mediana e bissetriz) em um triÂngulo equilátero de lado $l$.
	\color{black}
\end{exer}

%\subsection{Tabela Trigonométrica}

A partir dos valores encontrados no Exercício 1.2, construímos a seguinte tabela:
\begin{center}
	\begin{tabular}{|c|c|c|c|}
		\hline
		$\theta$ & $\sen\theta$ & $\cos\theta$ & $\tg\theta$ \\
		\hline
		$\pi/6$ & $1/2$ & $\sqrt{3}/2$ & $\sqrt{3}/3$ \\
		\hline
		$\pi/4$ & $\sqrt{2}/2$ & $\sqrt{2}/2$ & $1$ \\
		\hline
		$\pi/3$ & $\sqrt{3}/2$ & $1/2$ & $\sqrt{3}$ \\
		\hline
	\end{tabular}
\end{center}

A seguir resolveremos alguns exercícios que envolvem as razões métricas em um triângulo retângulo.
\begin{exer} \label{trigon1ex}
    Um avião decola, percorrendo uma trajetória retilínea, formando com o solo, um ângulo de $30^{\circ}$ (suponha que a região sobrevoada pelo avião seja plana). Depois de percorrer 1 000 metros, qual a altura atingida pelo avião? \\

\color{red}
A figura \ref{exercit1} descreve a situação
\begin{figure}[h!]
    \centering
    \includegraphics[scale=0.5]{exerc1t}
    \caption{Ilustração exercício \ref{trigon1ex}.}
    \label{exercit1}
\end{figure}
Note que temos envolvidos no problema um ângulo ($30^{\circ}$), o cateto oposto a esse ângulo ($h$) e a hipotenusa ($1000m$), assim, é conveniente usar o seno para obtermos a medida desconhecida.
$$\sin(30^{\circ})= \frac{h}{1000} \Longrightarrow \frac{1}{2}= \frac{h}{1000} \Longrightarrow h=500m$$
Assim, depois de percorrer $1000m$ o avião atinge $500m$.
\color{black}
\end{exer}

\begin{exer} \label{exerc34}
Um sensor detecta a aproximação de veículos em uma rua. Determine uma função $x(\theta)$ que forneça a distância do sensor ao veículo em relação a $\theta$. Se o sensor só detecta objetos quando $\theta$ é menor ou igual a $\pi/3$, determine a distância máxima que o carro pode estar do sensor para ser detectado. \\
 \begin{figure}[H]
	\centering
		\includegraphics[width=0.5 \textwidth,trim={3cm 13cm 3cm 15cm} ,clip]{figexercar}
  \caption{Figura para o exercício \ref{exerc34}.}
\end{figure}

\color{red}
A distância máxima ocorre quando o ângulo no qual o sensor detecta os objetos é máxima. Como, no exercício, temos o valor do cateto adjascente ao ângulo $\theta$ e queremos encontrar a hipotenusa. Portanto, iremos usar o cosseno.
$$\cos(60^{\circ})= \frac{8}{x} \Longrightarrow \frac{1}{2}= \frac{8}{x} \Longrightarrow x=16m$$
Portanto, a distância máxima que o carro pode estar do sensor para ser detectado é de $16m$.
\color{black}
\end{exer}

\section{Generalização da Noção de Ângulo}


Após uma volta completa, continuamos medindo os ângulos. Assim, por exemplo, o ângulo $\pi/6$ ocupa a mesma posição que os ângulos $2\pi+\pi/6$ (após uma volta) ou $4\pi+\pi/6$ (após duas voltas) e assim por diante. Mais geralmente, para qualquer $n$ natural, os ângulos $x$ e $x+2\pi n$ ocupam a mesma posição no círculo trigonométrico.

	Para dar sentido a ângulos negativos, convenciona-se que ao girarmos no sentido anti-horário, medimos os ângulos negativamente:

	\begin{figure}[h]
		\begin{center}
			\includegraphics[width=0.3\textwidth,trim={0 0 10,5cm 1cm},clip]{circ3}
		\end{center}
	\end{figure}

	Portanto, um ângulo $x$ radianos ocupa o mesmo lugar que todos os ângulos da forma
	\[x+2\pi n, ~n\in\Z.\]



\section{Generalização das Funções Trigonométricas: O círculo trigonométrico unitário.}

Considere o triângulo retângulo de hipotenusa 1, apresentado na figura \ref{trigon1}
  \begin{figure}[h!]
	\begin{center}
		\includegraphics[width=0.6 \textwidth,trim={3cm 13cm 3cm 15cm} ,clip]{figcirctrig1}
	\end{center}
 \caption{Triângulo retângulo de hipotenusa 1.}
\label{trigon1}
\end{figure}


Suponhamos que
\begin{itemize}
	\item o triângulo esteja situado no plano cartesiano,
	\item  o vértice associado ao ângulo $\theta$ esteja sobre a origem $(0,0)$,
	\item o cateto adjacente ao ângulo $\theta$ fique sobre a parte positiva do eixo $x$,
	\item  a hipotenusa do triângulo retângulo meça $1$ e o ângulo $\theta$ seja medido no sentido anti-horário a partir do eixo $x$.
\end{itemize}
Conforme observação anterior,

\begin{figure}[h]
	\begin{center}
		\includegraphics[width=0.6 \textwidth,trim={2cm 11cm 3cm 160mm} ,clip]{figcirctrig2}
	\end{center}
\end{figure}

Podemos variar os valores dos ângulos que para dois valores $\theta_1$ e $\theta_2$, teremos:

\begin{figure}[h!]
	\begin{center}
		\includegraphics[width=0.6 \textwidth,trim={3cm 165mm 2cm 6cm} ,clip]{figcirctrig3}
	\end{center}
	\label{trigon2}
\end{figure}

Podemos agora definir o círculo trigonométrico unitário (conforme figura \ref{trigon3}), que é o círculo do plano cartesiano de centro na origem $(0,0)$ e raio 1. O eixo $x$ é também chamado de \textit{eixo dos cossenos}, já o eixo $y$ é também chamado de \textit{eixo dos senos}.

Considere um ponto $A$ sobre o círculo trigonométrico no primeiro quadrante e seja $B$ a projeção desse ponto no eixo $x$, .  Seja $C$ a interseção do segmento de reta que passa por $\overline{OA}$ com a reta $x=1$ e $D=(1,0)$, ou seja, o segmento $\overline{CD}$ encontra-se sobre a reta $x=1$. Seja $\theta$ o ângulo $A\hat{O}B$. Observe que as coordendas de $A$ no plano cartesiano são $(\cos\theta,\sen\theta)$. Por semelhança de triângulos (triângulos $AOB$ e $COD$) vemos que
\[\frac{|\overline{CD}|}{1}=\frac{|\overline{CD}|}{|\overline{OD}|}=\frac{|\overline{AB}|}{|\overline{OB}|}=\frac{\sen\theta}{\cos\theta},\]
ou seja,
\[|\overline{CD}|=\tg\theta.\]

\begin{figure}[h!]
	\begin{center}
		\includegraphics[scale=3]{circtrigonom}
	\end{center}
	\caption{Círculo Trigonométrico Unitário.}
	\label{trigon3}
\end{figure}

Vemos assim que para $\theta\in(0,\pi/2)$ os valores $\sen\theta$, $\cos\theta$ e $\tg\theta$ podem ser obtidos a partir das coordenadas dos pontos $A$ e $C$: $A=(\cos\theta,\sen\theta)$ e $C=(1,\tg\theta)$. Mas essa construção pode ser feita para qualquer $\theta\in\R$ e dessa forma podemos definir as funções seno, cosseno e tangente para quaisquer ângulos reais.  Por exemplo, essa construção para um ângulo $\theta$ no segundo quadrante pode ser vista na figura \ref{trigon4}.

\begin{figure}[!h]\label{circtg}
	\begin{center}
		\includegraphics[width=0.4\textwidth,trim={0 1cm 9cm 3cm},clip]{circ5}
	\caption{Interpretação Geométrica da Tangente de um Ângulo no Segundo Quadrante}
	\end{center}
	\label{trigon4}
\end{figure}

Continuando a construção, podemos dividir o círculo trigonométrico unitário em quatro partes iguais, chamadas de quadrantes. No primeiro quadrante estão os ângulos entre $0$ e $90^{\circ}$, ou equivalentemente, entre $0$ e $\frac{\pi}{2}$. No segundo quadrante estão os ângulos entre $90^{\circ}$ e $180^{\circ}$ (ou entre $\frac{\pi}{2}$ e $\pi$). No terceiro quadrante estão os ângulos entre $180^{\circ}$ e $270^{\circ}$ (ou entre $\pi$ e $\frac{3\pi}{2}$). Por fim, no quarto quadrante estão os ângulos entre $270^{\circ}$ e $360^{\circ}$ (ou entre $\frac{3\pi}{2}$ e $2\pi$). Os ângulos $0,\frac{\pi}{2}, \pi, \frac{3\pi}{2}$ e $2\pi$ são denominados ângulos do eixo. Esses ângulos são notáveis e observando a figura \ref{trigon3}, podemos descobrir o valor do seno, cosseno e tangente desses arcos. A tabela \ref{tabelaeixos} apresenta esses valores.


\begin{table}[h!] \label{tabelaeixos}
\centering
\begin{tabular}{|l|c|c|c|c|r|}
	\hline
	 & $0$ & $\frac{\pi}{2}$  & $\pi$  & $\frac{3\pi}{2}$ & $2\pi$\\
	\hline
    $\sin(x)$ & 0 & 1 & 0& -1 & 0\\
	\hline
    $\cos(x)$ & 1 & 0 & -1 & 0 & 1\\
	\hline
	$\tan(x)$ & 0 & $\not\exists$ &0 & $\not\exists$ &0\\
	\hline
\end{tabular}
\caption{Seno, cosseno e tangente dos ângulos do eixo.}
\end{table}

\begin{obs} (\textbf{A relação fundamental trigonométrica})
A partir da figura \ref{trigon3}, por exemplo, podemos extrair uma relação trigonométrica, que é muito útil. Essa relação é conhecida como \textit{relação fundamental trigonométrica} e pode ser demonstrada, aplicando o Teorema de Pitágoras ao triângulo de hipotenusa 1 e catetos $\sin(\theta)$ e $\cos(\theta)$.
    $$\sin^2(x)+ \cos^2(x)=1$$
A partir dessa relação podemos obter duas outras conforme serão mostradas nos exercícios que seguem.
\end{obs}

 \begin{exer}
 Usando a relação fundamental trigonométrica mostre que : $\tan^2(x)+1= \sec^2(x)$, quando $\cos(x) \neq 0$ e $1+ \cot^2(x)= cossec^2(x)$ se $\sin(x) \neq 0$. \\
 \color{red}
 Da relação fundamental trigonométrica temos:
 \begin{eqnarray} \label{relafun}
 \sin^2(x)+\cos^2(x)=1
 \end{eqnarray}
 Sabendo que $\cos(x) \neq 0$, podemos dividir a equação \ref{relafun} por $\cos^2(x)$. Assim,

$$\frac{\sin^2(x)}{\cos^2(x)}+ 1= \frac{1}{\cos^2(x)}$$
Logo,
$$\tan^2(x)+1=\sec^2(x).$$
Para chegarmos na outra expressão basta dividirmos a relação fundamental por $\sin^2(x)$.

 \color{black}
 \end{exer}

\section{Fórmulas de redução ao primeiro quadrante}

Ao trabalharmos com Trigonometria e deparamo-nos com um ângulo que não se encontra no primeiro quadrante, sempre podemos reduzi-lo de forma a encontrar o ângulo correspondente a esse que esteja justamente no $1^{\circ}$ quadrante. Usamos esse procedimento em virtude ds simetria encontrada no ciclo trigonométrico. Mas precisamos nos atentar para o que ocorre com os sinais das funções trigonométricas em cada quadrante. A tabela \ref{sinal1} e a figura \ref{Sinai1} mostram os sinais em cada quadrante.
\begin{table}[H]
\centering
\begin{tabular}{|l|c|c|c|r|}
	\hline
	 & $1^{\circ} Q$ & $2^{\circ}$ Q & $3^{\circ}$ Q & $4^{\circ}$ Q\\
	\hline
    $\sin(x)$ & + & + & -& -\\
	\hline
    $\cos(x)$ & + & - & - & +\\
	\hline
	$\tan(x)$ & + & - &+ & -\\
	\hline
\end{tabular}
\label{sinal1}
\caption{Sinais em cada quadrante.}
\end{table}
\begin{figure}[H]
    \centering
    \includegraphics{sinais}
    \caption{Sinais em cada quadrante.}
    \label{Sinai1}
\end{figure}

\begin{itemize}
    \item Se o ângulo com o qual estamos trabalhando for $y$ e ele estiver no segundo quadrante, seu correspondente no $ 1^{\circ}$ quadrante será o ângulo $x$ tal que $\pi - x = y$ ou $180^{\circ} - x = y$. Você não deve se ater as fórmulas, a figura \ref{ciclo2p1} mostra que você pode fazer essa redução usando o círculo trigonométrico.
    \begin{figure}[H]
    \centering
    \includegraphics[scale=2]{ciclo2p1}
    \caption{Redução do $2^{\circ} $ quadrante para o $1^{\circ}$  quadrante.}
    \label{ciclo2p1}
\end{figure}
    \item Se o ângulo $y$ pertencer ao terceiro quadrante, seu correspondente $x$ no primeiro quadrante será dado por $x + \pi = y$ ou $180^{\circ} + x = y$. A figura \ref{ciclo2p2} mostra que você pode fazer essa redução usando o círculo trigonométrico.
     \begin{figure}[H]
    \centering
    \includegraphics[scale=2]{terp1}
    \caption{Redução do $3^{\circ} $ quadrante para o $1^{\circ}$  quadrante.}
    \label{ciclo2p2}
\end{figure}
    \item Se o ângulo com o qual estamos trabalhando for $y$ e ele estiver no quarto quadrante, seu correspondente no$ 1^{\circ}$ quadrante será o ângulo $x$ tal que $2\pi - x = y$ ou $360^{\circ} - x = y$. A figura \ref{ciclo2p3} mostra que você pode fazer essa redução usando o círculo trigonométrico.
     \begin{figure}[H]
    \centering
    \includegraphics[scale=2]{qp1}
    \caption{Redução do $4^{\circ} $ quadrante para o $1^{\circ}$  quadrante.}
    \label{ciclo2p3}
\end{figure}
\end{itemize}
\begin{obs}
    Tente n
\end{obs}
\begin{exer}
Reduza os ângulos $135^{\circ}$, $\frac{11 \pi}{6} rad$ e $\frac{4\pi}{3} rad$ e calcule os valores de $\sin(x)$, $\cos(x)$ e $\tan(x)$ para esses ângulos. \\
\color{red}
\begin{itemize}
    \item $x=135^{\circ}$ \\
    O arco $x=135^{\circ}$ está no $2^{\circ}$ quadrante, então para reduzirmos ao primeiro quadrante basta usarmos a fórmula $180^{\circ}-x=y$. Logo, $y=180^{\circ}- 135^{\circ}=45^{\circ}.$
     \begin{figure}[H]
    \centering
    \includegraphics[scale=2]{135}
    \caption{Redução do $135^{\circ} $ ao $1^{\circ}$  quadrante.}
    \label{ciclo2p1}
\end{figure}
Para encontrarmos as relações trigonométricas para esse ângulo, devemos nos atentar para o sinal das funções trigonométricas no quadrante a qual ele pertence. Nesse caso, temos que $x=135^{\circ}$ está no segundo quadrante. Assim,
\begin{eqnarray*}
\sin(135^{\circ})&=& \sin(45^{\circ})= \frac{\sqrt{2}}{2} \\
\cos(135^{\circ})&=& -\cos(45^{\circ})=- \frac{\sqrt{2}}{2} \\
\tan(135^{\circ})&=&- \tan(45^{\circ})= -1
\end{eqnarray*}
\item $x=\frac{11\pi}{6} rad$ \\
Sabendo que $\pi rad= 180^{\circ}$, então $\frac{11\pi}{6}= \frac{11.180}{6}=330^{\circ}$. Assim, $x=\frac{11\pi}{6} rad$ está no $4^{\circ}$ quadrante. Para reduzirmos ao primeiro quadrante, basta usarmos a fórmula $y=360^{\circ}-x$. Logo, $y=360^{\circ}-330^{\circ}=30^{\circ}$. Observando o sinal de cada função trigonométrica no $4^{\circ}$ temos:
\begin{eqnarray*}
\sin(330^{\circ})&=& -\sin(30^{\circ})= -\frac{1}{2} \\
\cos(330^{\circ})&=& \cos(30^{\circ})= \frac{\sqrt{3}}{2} \\
\tan(330^{\circ})&=& -\tan(30^{\circ})= -\frac{\sqrt{3}}{3}
\end{eqnarray*}
\item $x=\frac{4\pi}{3}$ \\
Sabendo que $x=\frac{4.180}{3}=240^{\circ}$, então $x \in 3^{\circ}$ quadrante e para reduzirmos ao terceiro quadrante usamos a relação $y=x-180^{\circ}$. Logo, $y=240^{\circ}-180^{\circ}=60^{\circ}$. Assim, observando o sinal de cada função trigonométrica no $3^{\circ}$ temos:
\begin{eqnarray*}
\sin(240^{\circ})&=& -\sin(60^{\circ})= -\frac{\sqrt{3}}{2} \\
\cos(240^{\circ})&=& \cos(60^{\circ})= \frac{1}{2} \\
\tan(240^{\circ})&=& -\tan(60^{\circ})= \sqrt{3}
\end{eqnarray*}
\end{itemize}
\color{black}
\end{exer}


\section{Gráficos de Funções Trigonométricas}

Um fenômeno é chamado periódico se o mesmo se repete sempre após o mesmo intervalo de tempo. Um exemplo simples de fenômeno periódico é o dia, determinado pelo movimento do Sol. A observação do movimento periódico dos corpos celestes nos possibilitou a criação do nosso calendário. Para uma melhor compreensão dos fenômemos periódicos, faz-se necessário o estudo das funções periódicas.

A repetição do valor numérico de uma função em um período determinado constitui uma função periódica. Mais precisamente, uma função $f$ é periódica de período $T$ se para todo $x$ real
\[f(x+T)=f(x).\]

Observe que conhecido o gráfico de uma função periódica em um determinado intervalo de seu domínio, que tenha o tamanho de seu período, automaticamente conhecemos  o gráfico em todo o domínio, uma vez que podemos repetir o gráfico em outros intervalos de mesmo tamanho. As funções trigonométricas são periódicas e por isso apresentaremos os seus gráficos em intervalos do tamanho de seus períodos. O período das funções seno e cosseno é $2\pi$ .  O período  da função tangente é $\pi$. \\
\begin{obs} As seguintes observações merecem destaque: \\
    \begin{itemize}
        \item As funções seno e cosseno são definidas em todo o conjunto dos números reais e seus valores variam entre $-1$ e $1$, ou seja, a imagem dessas funções é o intervalo fechado de $[-1,1]$. Isso significa que o valor máximo que a função pode atingir é $1$, enquanto o valor mínimo é $-1$.
        \item É importante notar que, além de possuir valor máximo de $1$ e valor mínimo de $-1$, as funções seno e cosseno atingem esses valores infinitas vezes ao longo de seus domínios.
        \item A função $y = \cos(x)$ é uma função par, o que significa que $\cos(-x) = \cos(x)$. Isso pode ser facilmente observado no círculo trigonométrico, já que a função cosseno representa a coordenada x de um ponto na circunferência unitária, que é simétrica em relação ao eixo y. Portanto, quando o ângulo é negativo, o cosseno mantém o mesmo valor que quando o ângulo é positivo.
        \item Para fazer o esboço do gráfico das funções $y=\sin(x)$ e $y=\cos(x)$, utilizamos os arcos do eixo no círculo trigonométrico.
        \item A função $y = \sin(x)$ é uma função ímpar, o que significa que $\sin(-x) = -\sin(x)$. Isso pode ser facilmente observado no círculo trigonométrico.
        \item Detalhes sobre função $y=\tan(x)$, será apresentado posteriormente, no entanto, vale adiantar que essa função é ímpar, pois $\tan(-x)= \frac{\sin(-x)}{\cos(-x)}= \frac{-\sin(x)}{\cos(x)}=-\tan(x)$.
    \end{itemize}
\end{obs}
As figuras \ref{senono} e \ref{cossenono} mostram o gráfico das funções $y=\sin(x)$ e $y=\cos(x)$ em um período, ou seja, no intevalo $[0,2\pi]$
\begin{figure}[h!]
	\begin{center}
		\includegraphics[scale=1.2]{seno-1}
	\end{center}
\caption{O gráfico da função seno em $[0,2\pi]$.}
\label{senono}
\end{figure}

\begin{figure}[h!]
	\centering
 \includegraphics[scale=1.2]{cos}
\caption{O gráfico da função cosseno em $[0,2\pi]$.}
\label{cossenono}
\end{figure}

 Os próximos exemplos apresentam como contruir gráficos de funções que são transformações nas funções $y=\sin(x)$ e $y=\cos(x)$.

 \begin{exer} Esboce o gráfico das funções seguintes, explicitando o domínio e a imagem.
\begin{itemize}
    \item[a)] $f(x)= 2 \sin(x) +2$ \\
    \color{red}
    Note que $D(f)= \mathbb{R}$, uma vez que não temos nenhuma restrição para $x$. \\
    Para encontrarmos a imagem de $f$ devemos lembrar que $-1 \leq \sin(x) \leq 1$ então se multiplicarmos por $2$ e depois somarmos $2$ em todos os termos da desigualdade chegamos que :
    $$-2+2 \leq 2 \sin(x) + 2 \leq 2+2 \Longrightarrow 0 \leq 2 \sin(x) + 2 \leq 4$$
Assim, $Im(f)=[0,4]$. \\
Uma vez que conhecemos o gráfico da função $y=\sin(x)$ podemos usar transformações para obtermos o gráfico de $f$. Uma vez que para obtermos a função $f$ multiplicamos a função $y=\sin(x)$ por $2$ e depois somamos $2$ unidades, então para fazermos o esboço do gráfico de $f$ precisamos fazer uma expansão vertical de comprimento $2$ e depois fazermos uma translação vertical (subimos $2$ unidades no eixo $y$).
\begin{figure}[h!]
    \centering
    \includegraphics[scale=0.9]{exemplo1seno}
    \caption{Em azul, o gráfico da função $y=2\sin(x)+2$ e em vermelho, gráfico da função $y=\sin(x)$.}
    \label{exercseno12}
\end{figure}
    \color{black}
    \item[b)] $g(x)= \frac{1}{2} \cos(x+ \frac{\pi}{2})-1$ \\

    \color{red}
        Note que $D(f)= \mathbb{R}$, uma vez que não temos nenhuma restrição para $x$. \\
    Para encontrarmos a imagem de $f$ devemos lembrar que $-1 \leq \cos\left(x+ \frac{\pi}{2} \right) \leq 1$ então se multiplicarmos por $\frac{1}{2}$ e depois subtraimos $1$ em todos os termos dessa desigualdade chegamos que :
    $$-1-\frac{1}{2}\leq \frac{1}{2} \cos(x+ \frac{\pi}{2}) -1 \leq \frac{1}{2} -1\Longrightarrow - \frac{3}{2} \leq \frac{1}{2} \cos(x+ \frac{\pi}{2}) - \leq -\frac{1}{2}$$
Assim, $Im(f)=[-\frac{3}{2},-\frac{1}{2}]$. \\
Conhecemos o gráfico da função $y=\cos(x)$ então para obtermos o gráfico da função $f$ precisamos fazer três transformações na função $y=\cos(x)$:
\begin{itemize}
    \item Contração vertical, por conta do fator $\frac{1}{2}$.
    \item deslocamento horizontal, ou seja, ''andar para trás'' $\frac{\pi}{2}$ unidades.
    \item Fazer uma translação vertical, ou seja, ''descer'' uma unidade no eixo $y$.
\end{itemize}
    \color{black}
\end{itemize}
\begin{figure}[h!]
    \centering
    \includegraphics[scale=0.9]{exerciciocosseno}
    \caption{Em azul, o gráfico da função $y=\frac{1}{2}\cos(x+\frac{\pi}{2})-1$ e em vermelho, gráfico da função $y=\cos(x)$.}
    \label{exercseno1}
\end{figure}
 \end{exer}
 \begin{obs}
     De modo geral, se temos uma função trigonométrica da forma $y=a+b \sin(kx+r)$ ou $y=a+b \cos(kx+r)$ a imagem de qualquer uma dessas funções é o intervalo $[a-b,a+b]$ e o período é $\frac{2\pi}{\vert k\vert}$. No entanto, ao invés de olharmos só para essas fórmulas é interessante entender o porquê disso.
 \end{obs}
\begin{obs}
	A função $f(x)=\tg x$ não está definida em $-\pi/2$ e $\pi/2$, e consequentemente (porque $f$ tem período $\pi$), em qualquer $x$ da forma $x=\pi/2+k\pi,$ para $k\in\mathbb{Z}$, logo, o domínio da função tangente é $\{x\in \mathbb{R};x\neq \pi/2+k\pi, \mbox{para todo }k\in\mathbb{Z}\}$ e a imagem é $\mathbb{R}$.
\end{obs}

\begin{figure}[h]
 \hspace{5cm}
 	\includegraphics[width=0.8\textwidth,trim={1cm 2cm 1cm 1cm},clip]{tg}
 \caption{O gráfico da função tangente em $(-\pi/2,\pi/2)$.}
\end{figure}



\section{Outras Funções Trigonométricas}

Definimos as funções seno, cosseno e tangente que são as principais funções trigonométricas, mas outras funções trigonométricas também são importantes:

\begin{description}
	\item \textbf{Cosecante de $x$:} $\ds\mbox{cossec} x=\frac{1}{\mbox{sen}x}$, se $\sin(x) \neq 0$.
	\item \textbf{Secante de $x$:} $\ds\mbox{sec} x=\frac{1}{\mbox{cos}x}$, se $\cos(x) \neq 0$.
	\item \textbf{Cotangente de $x$:} $\ds\mbox{cotg} x=\frac{1}{\mbox{tg}x}$, se $\tan(x) \neq 0$.
\end{description}

\begin{exer}
Pesquise sobre o domínio, a imagem, o gráfico e a inversa de cada uma das funções definidas nessa seção.
\end{exer}

\begin{exer} Se $\cos(x)= -\frac{2}{5} $ e $x \in 3^{\circ}$ quadrante, calcule o valor de $\sec(x)+\tan(x)$. \\
\color{red}
Sempre que tivermos o valor de uma das funções trigonométricas $\sen(x)$ ou $\cos(x)$, a \textit{Relação Fundamental Trigonométrica} será útil para encontrar a outra função. Por exemplo, nesse exercício temos o valor de $\cos(x)$, então podemos encontrar o valor de $\sin(x)$ usando a \textit{Relação Fundamental Trigonométrica}. Assim,
\begin{eqnarray*}
\sin^2(x)+\cos^2(x)&=& 1 \\
\sin^2(x)+\left( -\frac{2}{5} \right)^2&=&1 \\
\sin^2(x)+ \frac{4}{25}&=& 1 \\
\sin^2(x)&=& 1- \frac{4}{25} \\
\sin^2(x) &=& \frac{21}{25} \\
\sin(x)&=&  \pm \frac{\sqrt{21}}{5}
\end{eqnarray*}
Uma vez que $x \in 3^{\circ}$ quadrante, então $\sin(x)=- \frac{\sqrt{21}}{5}$, assim:
$$\sec(x)= \frac{1}{\cos(x)}=-\frac{5}{2}$$
e
$$\tan(x)= \frac{\sin(x)}{\cos(x)}=\frac{-\frac{\sqrt{21}}{5}}{-\frac{2}{5}}=\frac{\sqrt{21}}{2}$$
Consequentemente
$$\sec(x)+\tan(x)= -\frac{5}{2}+\frac{\sqrt{21}}{2}= \frac{-5+\sqrt{21}}{2}.$$
\color{black}
\end{exer}


\begin{exer} Se $\sin(x) = \frac{3}{5}$ e $x \in 1^{\circ}$ quadrante, encontre o valor da expressão
$$z=\frac{\sec(x)- cossec(x)}{\tan(x)+ \ds\mbox{cotg} (x)}$$
\color{red}
Usando a \textit{Relação Fundamental Trigonométrica} temos :
\begin{eqnarray*}
\sin^2(x)+ \cos^2(x)&=&1 \\
\left( \frac{3}{5}\right)^2 + \cos^2(x)&=&1 \\
\cos^2(x) &=& 1-\frac{9}{25} \\
\cos^2(x) &=& \frac{16}{25} \\
\cos(x) &=& \pm \frac{4}{5}
\end{eqnarray*}
Sendo $x \in 1^{\circ}$ quadrante, então $\cos(x)= \frac{4}{5}$, assim:
\begin{itemize}
    \item $\sec(x)=\frac{1}{\cos(x)}=\frac{5}{4}$ ;
    \item $\tan(x)=\frac{\sin(x)}{\cos(x)}= \frac{\frac{3}{5}}{\frac{4}{5}}= \frac{3}{4}$;
    \item $cossec(x)= \frac{1}{\sin(x)}= \frac{5}{3}$;
    \item $cotg(x)= \frac{1}{\tan(x)}= \frac{4}{3}$.
\end{itemize}
Portanto,
$$z=\frac{\sec(x)- cossec(x)}{\tan(x)+ \ds\mbox{cotg} (x)}= \frac{\frac{5}{4}-\frac{5}{3}}{\frac{3}{4}+\frac{4}{3}}= -\frac{1}{5}.$$
\color{black}
\end{exer}
\section{Identidades Trigonométricas}

Algumas relações conhecidas como \textbf{Identidades Trigonométricas} serão fundamentais no curso:

\begin{enumerate}
	\item (\textbf{Seno da soma}) $\sen(a+b)=\sen a\cos b+\sen b\cos a$.
	\item (\textbf{Cosseno da soma:}) $\cos(a+b)=\cos a\cos b-\sen b\sen a$. \\
 \textbf{Para visualizar as demonstrações dos itens 1 e 2, assista ao vídeo disponível no seguinte link:}  \\
 \begin{center}
     \url{https://www.youtube.com/watch?v=MlCZU4zNymo}
 \end{center}
	\item (\textbf{Tangente da soma})$\tg(a+b)=\ds\frac{\tg a+\tg b}{1-\tg a\tg b}$.
 \begin{proof}
 \color{red}
 Para demonstrar esse item, usamos a definição de função tangente e as expressões dos itens 1 e 2. Assim,
 $$\tg(a+b)=\frac{\sen(a+b)}{\cos(a+b)}=\frac{\sen a\cos b+\sen b\cos a}{\cos a\cos b-\sen b\sen a}$$
 Dividindo o numerador e o denominador por $\cos a\cos b$, desde que seja não nulo, temos:
 $$\tg(a+b)=\frac{\frac{\sen a\cos b+\sen b\cos a}{\cos a\cos b}}{\frac{\cos a\cos b-\sen b\sen a}{\cos a\cos b}}=\ds\frac{\tg a+\tg b}{1-\tg a\tg b}$$
 Portanto, a fórmula para a tangente da soma é:

$$\boxed{\tg(a+b)=\ds\frac{\tg a+\tg b}{1-\tg a\tg b}}$$
 \color{black}
 \end{proof}
 \item (\textbf{Seno da diferença:}) $\sen(a-b)=\sen a\cos b-\sen b\cos a$.
 \item (\textbf{Cosseno da diferença:}) $\cos(a-b)=\cos a\cos b+\sen b\sen a$. \\
 \textbf{O itens 4 e 5 podem ser demonstrados trocando b por $-b$ nas expressões 1 e 2, respectivamente e usando o fato de que a funação seno é ímpar e a função cosseno é par.} \\
 \item (\textbf{Arco duplo: Seno }) $\sin(2a) = 2\sin(a)\cos(a)$
 \begin{proof}
 \color{red}
     Essa fórmula é muito útil na trigonometria, pois permite calcular o valor do seno de um ângulo duplo em termos dos valores de seno e cosseno do ângulo original. E para demonstrá-la, vamos partir da fórmula da adição para o seno, apresentada no item 1, trocando $b$ por $a$. Logo, \\
     $$\sin(2a) = \sin(a + a) = \sin(a)\cos(a) + \cos(a)\sin(a)=2\sin(a)\cos(a)$$


Portanto, a fórmula para o seno do arco duplo é:

$$\boxed{\sin(2a) = 2\sin(a)\cos(a)}$$
     \color{black}
 \end{proof}
 \item (\textbf{Arco duplo: Cosseno }) $\cos(2a)=\cos^2(a)-\sin^2(a)$.
 \begin{proof}
 \color{red}
     Vamos partir da fórmula da adição para o cosseno, apresentada no item 2:

$$\cos(\alpha + \beta) = \cos(\alpha)\cos(\beta) - \sin(\alpha)\sin(\beta)$$

Substituindo $\alpha = a$ e $\beta = a$, temos:

$$\cos(2a) = \cos(a + a) = \cos(a)\cos(a) - \sin(a)\sin(a)= \cos^2(a)-\sin^2(a).$$
Portanto, a fórmula para o seno do arco duplo é:
\begin{eqnarray}
\boxed{\cos(2a) = cos^2(a)-\sin^2(a)}
\label{duplo}
\end{eqnarray}
 \end{proof}
 Note isolando $\sin^2(x)$ na identidade trigonométrica $\sin^2(a) + \cos^2(a) = 1$, podemos reescrever a expressão \label{duplo} como:

$$\cos(2a) = \cos^2(a) - (1 - \cos^2(a))$$

Simplificando, obtemos uma outra expressão para $\cos(2a)$, que depende apenas de cosseno:
\begin{eqnarray}
\cos(2a) = 2\cos^2(a) - 1
\label{duplo1}
\end{eqnarray}

Se isolarmos $\cos^2(x)$ na identidade trigonométrica $\sin^2(a) + \cos^2(a) = 1$, podemos reescrever a expressão \label{duplo} como:
\begin{eqnarray}
\cos(2a) = 1-2\sin^2(a)
\label{duplo2}
\end{eqnarray}
Essas expressões serão essenciais para demonstrar as fórmulas de arco metade que seguem:
\item  (\textbf{Arco metade: Cosseno }) $\cos^2(x)= \frac{1}{2}(1+\cos(2x)) $
\begin{proof}
    \color{red}
    Começamos com a fórmula \ref{duplo1},  $\cos(2a) = 2\cos^2(a) - 1$ e substituímos $a$ por $\frac{x}{2}$ para obter:

$$\cos(x) = 2\cos^2\left(\frac{x}{2}\right) - 1$$

Isolando $\cos^2\left(\frac{x}{2}\right)$, temos:

$$\cos^2\left(\frac{x}{2}\right) = \frac{1}{2}(1 + \cos(x))$$

Substituindo $\frac{x}{2}$ por $x$ e reescrevendo a equação em termos de $\cos(2x)$, temos:

$$\cos^2(x) = \frac{1}{2}(1 + \cos(2x))$$

Portanto, $$\boxed{\cos^2(x)= \frac{1}{2}(1+\cos(2x))}$$

    \color{black}
\end{proof}
\item  (\textbf{Arco metade: Seno }) $\sin^2(x)= \frac{1}{2}(1-\cos(2x)) $
\begin{proof}
    \color{red}
    Para comprovar esse item, é possível utilizar a mesma abordagem utilizada anteriormente, porém aplicando agora a fórmula \ref{duplo2}. Dessa forma, será possível demonstrar de maneira equivalente a veracidade desse item.
    \color{black}
\end{proof}
\end{enumerate}
\section{Equações Trigonométricas}
Uma equação trigonométrica é uma equação que envolve funções trigonométricas como seno, cosseno, tangente, cotangente, secante ou cossecante de uma variável trigonométrica. O objetivo da resolução de equações trigonométricas é encontrar todas as soluções possíveis para a equação.

As equações abaixo são exemplos de equações trigonométricas:
\begin{enumerate}
    \item $\sin(x)= \frac{1}{2}$
    \item $\cos^2(x)-\cos(x)=0$
    \item $\cos(x)(\tan(x)-1)=0$
\end{enumerate}

Para resolvermos uma equação trigonométrica, precisamos nos atentar em qual intervalo é pedida a solução.

A resolução de equações trigonométricas pode ser feita por meio de diversas técnicas, dependendo da complexidade da equação. Alguns dos métodos mais comuns incluem:
\begin{enumerate}
    \item Equações simples: são aquelas que envolvem apenas uma função trigonométrica. Nesse caso, é possível usar propriedades básicas das funções trigonométricas para encontrar as soluções.

 \item Equações quadráticas: são aquelas que envolvem o quadrado de uma função trigonométrica. Nesse caso, é possível fazer uma substituição adequada para transformar a equação em uma equação quadrática comum, que pode ser resolvida usando as técnicas usuais.

\item Equações com produto nulo: são aquelas em que uma ou mais funções trigonométricas multiplicadas entre si resultam em zero. Nesse caso, é possível usar a propriedade do produto nulo para encontrar as soluções.

\item Uso de identidades: muitas vezes, as equações trigonométricas podem ser resolvidas por meio da aplicação de identidades trigonométricas, como as identidades de soma e diferença, identidades de produto e identidades de duplo ângulo.
\end{enumerate}
Em geral, a resolução de equações trigonométricas requer um bom conhecimento das propriedades e identidades das funções trigonométricas, bem como das técnicas algébricas usadas para resolver equações.
 Vamos resolver alguns exemplos:

 \begin{exe}
     Resolva a seguinte equação trigonométrica $\sin(\theta) =\frac{\sqrt{2}}{2}$:
     \begin{itemize}
         \item[a)] No intervalo $[0,2\pi]$: \\
         \color{blue}
    Note que, nesse caso, queremos encontrar os arcos no intervalo  $[0,2\pi]$, cujo seno vale $\frac{\sqrt{2}}{2}$. Sabemos que $\sin(\theta)$ é positivo no primeiro e segundo quadrante. No primeiro quadrante o ângulo é $\frac{\pi}{4}$ (usamos sempre o arco do $1^{\circ}$ como referência). Para encontrar o arco no segundo quadrante, usamos a expressão $\pi-\frac{\pi}{4}= \frac{3\pi}{4}$, conforme mostra a figura  \ref{equa1}.
         \begin{figure}[h!]
	\centering
 \includegraphics[scale=2]{equa1png}
\caption{Ilustração da sulução da equação $\sin(\theta) =\frac{\sqrt{2}}{2}$.}
\label{equa1}
\end{figure}
Portanto, a solução da equação pedida é $S= \{ \frac{\pi}{4}, \frac{3\pi}{4} \}$ .
         \color{black}
         \item[b)] Encontre a solução geral: \\
         \color{blue}
         A solução geral é dada por:
         $$S= \boxed{ \{ x \in \mathbb{R} | x=\frac{\pi}{4}+2k\pi \,\, ou \,\, x= \frac{3\pi}{4} +2k\pi, \,\, k \in \mathbb{Z} \}}$$.
         \color{black}
     \end{itemize}
 \end{exe}
 \begin{exe}
     Resolva a equação $2\cos(\frac{\theta}{3}) -
\sqrt{3} = 0,$ onde $0 \leq \theta \leq 3\pi$. \\
\color{blue}
A equação  $2\cos(\frac{\theta}{3}) -
\sqrt{3} = 0$ é equivalente a  $\cos(\frac{\theta}{3}) =
\frac{\sqrt{3}}{2}$. Como $0 \leq \theta \leq 3\pi$, então, $0 \leq \frac{\theta}{3} \leq \pi $. Uma vez que o único arco no intervalo $[0,\pi]$, que possui cosseno igual a $\frac{\sqrt{3}}{2}$ é  $\frac{\pi}{6}$, temos que:
$$\frac{\theta}{3}= \frac{\pi}{6} \Longrightarrow \theta=\frac{\pi}{2}$$. Portanto,
$$ \boxed{S= \{\frac{\pi}{2} \}}$$.
\color{black}
 \end{exe}
\begin{exe} Resolva a equação $2\cos(x)= -1$ no intervalo $[0,2\pi]$. \\
\color{blue}
Primeiro, vamos isolar o $\cos(x)$ dividindo ambos os lados por 2:
$$\cos(x)= -\frac{1}{2}$$

Sabemos que o cosseno é negativo no segundo e terceiro quadrante, para encontrarmos quais são esses arcos, pegamos como referência o arco no primeiro quadrante que tem cosseno igual a $\frac{1}{2}$, isto é, $\frac{\pi}{3}$  . Conforme mostra a figura \ref{equa2}, o arco no segundo quadrante é $\pi-\frac{\pi}{3}= \frac{2\pi}{3}$ e o arco no terceiro quadrante é $\pi+\frac{\pi}{3}= \frac{4\pi}{3}$. Assim, $\cos(x)=-\frac{1}{2}$ ocorre quando $x$ é um ângulo de $\frac{2\pi}{3}$ ou $\frac{4\pi}{3}$ (ou seus equivalentes negativos). Portanto, a solução da equação no intervalo $[0,2\pi]$ é:
$$S=\boxed{ \{\frac{2\pi}{3}, \frac{4\pi}{3} \}}$$


 \begin{figure}[h!]
	\centering
 \includegraphics[scale=2]{equa2}
\caption{Ilustração da solução da equação $2\cos(x)= -1$ .}
\label{equa2}
\end{figure}
\color{black}
\end{exe}

\begin{exe}
    Resolva a equação $(\sin(\theta) -\frac{1}{2}) (\tan(\theta) -1) = 0$, para $0 \leq \theta  <  \frac{\pi}{2}$. \\
    \color{red}
    Essa equação é do tipo produto nulo. Sabemos que produto $ab$ só é zero se $a = 0$ ou $ b = 0$. Portanto:
    $$\sin(\theta) -\frac{1}{2}=0 \,\, ou \,\,\tan(\theta) -1=0 $$
    Analisando cada caso, separado:
    $$\sin(\theta) -\frac{1}{2}=0 \Longrightarrow \sin(\theta)=\frac{1}{2} \Longrightarrow \theta= \frac{\pi}{6}$$
    $$\tan(\theta) -1=0 \Longrightarrow \tan(\theta) =1 \Longrightarrow \theta=\frac{\pi}{4}$$
     Assim, a solução da equação pedida é:
     $$\boxed{S=\{\frac{\pi}{6}, \frac{\pi}{4}  \}}.$$
    \color{blue}
\end{exe}

\begin{exe}
    Resolva a equação  $2\sin^2 x - \sin x - 1 = 0$ no intervalo $[0,2\pi]$. \\
    \color{red}
   Neste caso, temos uma equação quadrática trigonométrica que pode ser resolvida através da mudança de variável $y = \sin(x)$. Dessa forma, obtemos a equação do segundo grau $2y^2  - y - 1 = 0$, cujas raízes são $y = 1$ e $y = -\frac{1}{2}$. Portanto, podemos encontrar os valores de $x$ correspondentes da seguinte maneira:

Para $y = 1$, temos $\sin(x) = 1$, o que implica em $x = \frac{\pi}{2}$.

Para $y = -\frac{1}{2}$, temos $\sin(x) = -\frac{1}{2}$. Isso ocorre em dois quadrantes, $3^{\circ}$ e $4^{\circ}$ quadrantes. Portanto, os valores possíveis de $x$ são $\pi + \frac{\pi}{6} = \frac{7\pi}{6} (3^{\circ}Q)$ ou $2\pi - \frac{\pi}{6} = \frac{11\pi}{6} (4^{\circ}Q)$.
    Assim o conjunto solução da equação original é:
    $$\boxed{S=\{\frac{\pi}{2},\frac{7\pi}{6}, \frac{11\pi}{6}\}}$$

    \color{black}
\end{exe}
\section{Funções Trigonométricas Inversas}

As funções seno, cosseno e tangente não são bijetoras, mas uma restrição no domínio das mesmas nos possibilita definir suas inversas. Essa escolha de domínio não é única, mas adotaremos as escolhas padrões, que são os intervalos que contém os ângulos notáveis e cujas imagens das funções restritas são iguais às imagens das funções sem restrição. Os gráficos das funções inversas são obtidos a partir dos gráficos das funções originais através da reflexão do gráfico da função original em torno da reta $y=x$. Por isso apresentaremos três imagens: o gráfico das funções trigonométricas restritas, esse gráfico juntamente com sua reflexão e o gráfico das funções trigonométricas inversas.

\subsection{Arco Seno}

A função
\[\sen: [-\pi/2,\pi/2]\to[-1,1]\]
é bijetora e sua inversa é denominada arco seno: \[\arcsen:[-1,1]\to[-\pi/2,\pi/2].\] Por definição,
\[\arcsen x=y \Longleftrightarrow \sen y=x.\]

Por serem inversas uma da outra, essas duas funções satisfazem:
\begin{description}
	\item[1:] $\sen(\arcsen(x))=x,\quad\forall\quad x ~\in [-1,1].$
	\item[2:] $\arcsen(\sen(x))=x,\quad \forall\quad x~ \in [-\pi/2,\pi/2]$.
\end{description}

Para encontrarmos o gráfico da função arcoseno, vamos analisar o  gráfico da função seno em $[-\pi/2,\pi/2]$ como dado a seguir:

\begin{figure}[h!]
	\begin{center}
		\includegraphics[width=0.6\textwidth,trim={1cm 10cm 12cm 3cm},clip]{sen1}
	\end{center}
\end{figure}

A reflexão do gráfico de seno em torno no eixo $x$:

\begin{figure}[h!]
	\begin{center}
		\includegraphics[width=0.6\textwidth,trim={1cm 4cm 12cm 2cm},clip]{sen2}
	\end{center}
\end{figure}

E finalmente, o gráfico do arco seno:

\begin{figure}[h!]
	\begin{center}
		\includegraphics[width=0.6\textwidth,trim={1cm 6cm 12cm 3cm},clip]{sen4}
	\end{center}
\end{figure}

\begin{exer}
	Calcule o valor de arcsen$x$ para $x$ assumindo os seguintes valores: $-\frac{\sqrt{3}}{2}, 0, \frac{\sqrt{2}}{2}, \frac{1}{2}, 1$  \\
	\color{red}
	Temos que:
\begin{itemize}
    \item $\arcsin\left(-\frac{\sqrt{3}}{2}\right)=-\frac{\pi}{3}$.
\item $\arcsin(0)=0$.
\item $\arcsin\left(\frac{\sqrt{2}}{2}\right)=\frac{\pi}{4}$.
\item $\arcsin\left(\frac{1}{2}\right)=\frac{\pi}{6}$.
\item $\arcsin(1)=\frac{\pi}{2}$.
\end{itemize}
Vale lembrar que o valor de $\arcsin(x)$ está restrito ao intervalo $\left[-\frac{\pi}{2},\frac{\pi}{2}\right]$. Por exemplo, no caso de $\arcsin\left(-\frac{\sqrt{3}}{2}\right)$, o valor $-\frac{\pi}{3}$ está nesse intervalo, mas o valor $\frac{2\pi}{3}$ também é uma solução da equação $\sin(x)=-\frac{\sqrt{3}}{2}$, porém não está no intervalo $\left[-\frac{\pi}{2},\frac{\pi}{2}\right]$. Por isso, é importante sempre verificar se a solução encontrada está dentro do intervalo de definição de $\arcsin(x)$.
	\color{black}
\end{exer}

\subsection{Arco Cosseno}

A função
\[\cos: [0,\pi]\to[-1,1]\]
é bijetora e sua inversa é denominada arco cosseno:
\[\arccos:[-1,1]\to[0,\pi].\]

Por definição,
\[\arccos x=y \Longleftrightarrow \cos y=x.\]


Por serem inversas uma da outra, essas duas funções satisfazem:
\begin{description}
	\item[1:] $\cos(\arccos(x))=x,\quad\forall\quad x ~\in [-1,1].$
	\item[2:] $\arccos(\cos(x))=x,\quad \forall\quad x~ \in [0,\pi]$.
\end{description}

A seguir o gráfico da função cosseno em $[0,\pi]$ e de sua reflexão em torno de $y=x$:

\begin{figure}[h]
	\begin{center}
		\includegraphics[width=0.4\textwidth,trim={1cm 4cm 12cm 2cm},clip]{cos2}
		\includegraphics[width=0.4\textwidth,trim={1cm 3cm 12cm 1cm},clip]{cos3}
	\end{center}
\end{figure}

\newpage

O gráfico do arco cosseno:

\begin{figure}[h]
	\begin{center}
		\includegraphics[width=0.5\textwidth,trim={1cm 6cm 20cm 4cm},clip]{cos4}
	\end{center}
\end{figure}


\begin{exer}
	Calcule o valor de arccos$x$ para $x$ assumindo os seguintes valores: $-\frac{\sqrt{3}}{2}, 0, \frac{\sqrt{2}}{2}, \frac{1}{2}, 1$  \\
\end{exer}
\color{red}
Temos que:
\begin{itemize}
\item $\arccos\left(-\frac{\sqrt{3}}{2}\right)=\frac{5\pi}{6}$.
\item $\arccos(0)=\frac{\pi}{2}$.
\item $\arccos\left(\frac{\sqrt{2}}{2}\right)=\frac{\pi}{4}$.
\item $\arccos\left(\frac{1}{2}\right)=\frac{\pi}{3}$.
\item $\arccos(1)=0$.
\end{itemize}
Vale lembrar que o valor de $\arccos(x)$ está restrito ao intervalo $[0,\pi]$. Por exemplo, no caso de $\arccos\left(-\frac{\sqrt{3}}{2}\right)$, o valor $\frac{7\pi}{6}$ é uma solução da equação $\cos(x)=-\frac{\sqrt{3}}{2}$, porém não está no intervalo $[0,\pi]$. Por isso, é importante sempre verificar se a solução encontrada está dentro do intervalo de definição de $\arccos(x)$.
\color{black}
\subsection{Arco Tangente}

A função
\[\tg: (-\pi/2,\pi/2)\to \R\]
é bijetora e sua inversa é denominada arco tangente:
\[\arctg:\R\to(-\pi/2,\pi/2).\]

Por definição,
\[\arctg x=y \Longleftrightarrow \tg y=x.\]

Por serem inversas uma da outra, essas duas funções satisfazem:
\begin{description}
	\item[1:] $\tg(\arctg(x))=x,\quad\forall\quad x ~\in \R.$
	\item[2:] $\arctg(\tg(x))=x,\quad \forall\quad x~ \in (-\pi/2,\pi/2)$.
\end{description}

A seguir o gráfico da função tangente em $(-\pi/2,\pi/2)$ e de sua reflexão em torno de $y=x$:

\begin{figure}[h]
	\begin{center}
		\includegraphics[width=0.4\textwidth,trim={1cm 2cm 12cm 2cm},clip]{tg1} \hspace{1cm}
		\includegraphics[width=0.4\textwidth,trim={1cm 1cm 10cm 1cm},clip]{tg2}
	\end{center}
\end{figure}

O gráfico do arco tangente:

\begin{figure}[h]
	\begin{center}
		\includegraphics[width=0.6\textwidth,trim={1cm 6cm 10cm 3cm},clip]{tg4}
	\end{center}
\end{figure}


\begin{exer}
	Calcule o valor de arctg$x$ para $x$ assumindo os seguintes valores: $-\frac{\sqrt{3}}{3}, 0, \sqrt{3}, 1$  \\
\end{exer}
\color{red}
Temos que:
\begin{itemize}
    \item $\arctan\left(-\frac{\sqrt{3}}{3}\right)=-\frac{\pi}{6}$.
\item $\arctan\left(0\right)=0$.
\item $\arctan\left(\sqrt{3}\right)= \frac{\pi}{3}$
\item $\arctan\left(1\right)=\frac{\pi}{4}$
\end{itemize}
Vale lembrar que o valor de $\arctan(x)$ está restrito ao intervalo $\left(-\frac{\pi}{2},\frac{\pi}{2}\right)$. Por exemplo, no caso de $\arctan\left(-\frac{\sqrt{3}}{3}\right)$, a solução encontrada $-\frac{\pi}{6}$ está dentro do intervalo $\left(-\frac{\pi}{2},\frac{\pi}{2}\right)$, porém é importante sempre verificar se a solução encontrada está dentro do intervalo de definição de $\arctan(x)$.
\color{black}
\section{Exercícios}
\begin{enumerate}
\item A polícia federal localizou na floresta amazônica uma pista de pouso clandestina com as
seguintes características:
\begin{itemize}
\item A pista media 300 metros de comprimento, era plana e horizontal;
\item No final da pista havia uma árvore de 30 m de altura.
\end{itemize}
Um pequeno avião encontra-se em um ponto A da pista a 300 metros da árvore. O avião partirá desse ponto A, em direção à árvore e levantará vôo após percorrer 180 metros de pista. Se o avião levantar vôo em linha reta, de modo que essa reta (trajetória do avião no ar) forme um ângulo $\alpha$ com o plano horizontal, qual deverá ser a tangente de $\alpha$ para que a aeronave passe exatamente a 10 metros acima da árvore?

\item Uma TV encontra-se fixada em uma parede e a dois metros de distância dessa parede está posicionado um sofá de três lugares. Uma pessoa senta-se em um dos lugares das extremidades dos sofá de forma que sua distância até a TV é de 3 metros. O ângulo de visão dessa pessoa é o ângulo formado entre a reta que liga a visão da pessoa ao centro da TV e a reta perpendicular à parede passando pelo centro da TV. Sabendo que a visão da TV só é boa se a pessoa assiste TV em um ângulo menor ou igual a $\pi/3$, verifique se essa pessoa terá uma boa visão da TV.
\item Dê os valores dos ângulos a seguir e radianos e calcule o valor das funções seno, cosseno, tangente, secante, cossecante e cotangente nos mesmos:
\begin{enumerate}
	\item $150^{\circ}$
	\item $315^{\circ}$
	\item $210^{\circ}$
\end{enumerate}

\item Sabendo que $\sen(x)=-\frac{3}{4}$ e $\pi \leq x \leq \frac{3\pi}{2}$, determine o valor de $\cos(x)$, $\tg(x)$ e $\sec(x)$

 \item Sabendo que $\sin(\theta)=\frac{1}{4}$ e que $\theta$ está no segundo quadrante, determine o valor de $\cos(\theta)$, $\tan(\theta)$ e $\sec(\theta)$.
 \item Sabendo que $\cos(x)= \frac{3}{5}$, determine os possíveis valores de $\sin(x)$ e os quadrantes aos quais $x$ pode pertencer.
\item Trace os gráficos das funções a seguir incluindo ao menos um período, indique também qual a imagem, o valor máximo e o valor mínimo de cada função:
\begin{enumerate}
\item  $f(x)=\mbox{tg}(x+\pi/4)$
\item $f(x)=\mbox{cos}(4x)$
\item $f(x)=1-2\sin (x)$
\item $f(x)=\sen(x)+2$
\end{enumerate}

\item Simplifique as expressões trigonométricas, supondo sempre que os denominadores que aparecer são não nulos.
\begin{itemize}
    \item[a)] $\ds\frac{\mbox{sen}^2x}{\mbox{cos}x}-\mbox{sec}x$.
    \item[b)] $\frac{\sen(4x)}{\sen(x)\cos(2x)}$
    \item[c)] $\mbox{cotg}x \hspace{0,1cm} \mbox{sec}x \hspace{0,1cm} \mbox{cossec}x $
\end{itemize}
 \item Faça a redução ao  primeiro quadrante e encontre o valor de:
     \begin{enumerate}
         \item $cossec(225^{\circ})$
         \item $cotg(150^{\circ})$
         \item $tg\left(\frac{3\pi}{4})- cotg(\frac{3\pi}{4}\right)$
         \item $\cos\left(-\frac{2\pi}{3}\right)$
         \item $\sin(630^{\circ})$
         \item $\cos\left( \frac{37\pi}{6}\right)- \sin(6\pi)$
         \item $\tan(315^{\circ})-\cos(330^{\circ})$
     \end{enumerate}

\item  Qual o maior e o menor valor que $f(x) = 5+5\sen(5x+5)$ assume?

\item A altura da cabine de uma roda gigante é descrita em função do tempo (em min) por
$$h(t)=76+75 \sin \left( \frac{\pi t}{15}- \frac{\pi}{2}\right)$$
\begin{enumerate}
    \item Determine as alturas máxima e mínima da cabine.
    \item Determine quanto tempo a roda gigante demora para dar uma volta completa (ou seja, o período de $h(t)$).
    \item Usando um programa gráfico, trace a curva $h(t)$.
\end{enumerate}
\item Usando a fórmula de adição de arcos, encontre os valores de $\sen(75^{ \circ})$, $\cos(75^{ \circ})$ e $\tg(75^{ \circ})$. Dica: Note que $75^{\circ}=30^{\circ}+ 45^{\circ}$.
 \item Nesse exercício você poderá usar as seguintes relações trigonométricas:
    $$\sin(x+y)= \sin(x)\cos(y)+\sin(y)\cos(x) \,\, e \,\, \cos(x+y)=\cos(x)\cos(y)-\sin(x)\sin(y)$$
    \begin{enumerate}
        \item Encontre uma expressão para $\cos(3x)$.
        \item  Mostre que podemos escrever $\cos(2x)=1-2\sin^2(x)$.
        \item Mostre que $\tan(2x)=\frac{2\tan(x)}{1-\tan^2(x)}$
        \item Determine $\sin(x-y)$, $\cos(x-y)$ e $\tan(x-y)$, sabendo que $\sin(x)=\frac{4}{5}$, $\cos(y)=\frac{\sqrt{5}}{5}$ e que $x,y \in 1^{\circ}$ quadrante.
        \item  Encontre o valor de $\sin(15^{\circ})$ e a partir desse resultado, encontre o valor de $\sin(7,5^{\circ})$.
        \item Mostre que $\cos \left(x + \frac{\pi}{2}\right)=-\sin(x)$.
        \item Encontre o valor da expressão $\ln(\sec(\theta))+\ln(\cos(\theta))$
        \item Econtre o valor da espressão
        $$\cos(\arctan(x)+\arcsin(\frac{1}{\sqrt{3}}))$$
        obs: Considere que os ângulos envolvidos nesse item estão sempre no primeiro quadrante.
    \end{enumerate}
    \item Sabe-se que $x + y = \frac{\pi}{4}$ e $\sin(x) = \frac{5}{13}$, com $0 \leq x \leq \frac{\pi}{2}$, calcule $\sin(y)$ e $\cos(y)$.
\item Simplifique a expressão abaixo, levando em conta que os denominadores são sempre não nulos.
$$\frac{\sin(\frac{\pi}{2}-x)\sin(\pi-x)}{\cos(\pi-x)\cos(2\pi-x)}$$
\item Se $\sin(x) = m$ e $\cos(x) = n$, determine o valorde $\sin(2x)$, $\cos(2x)$ e $\tan(2x)$.
\item Calcule o valor de $\sin(105^{\circ})$, $\sec(105^{\circ})$, $\cos(15^{\circ})$ e $\tan(15^{\circ})$.
\item Se $\sin(x)= \frac{1}{3}$ e $x \in 2^{\circ}$ quadrante, encontre o valor de $\sin(\frac{x}{2})$, $\cos(\frac{x}{2})$ e $\tan(\frac{x}{2})$
\item A partir do gráfico de $f(x)= \cos(x)$, faça o esboço do gráfico de $g(x)=-1+ 2 \cos(x+ \frac{\pi}{2})$. Encontre também o período e a imagem.


\item Resolva as seguintes equações trigonométricas, no intervalo especificado:
\begin{enumerate}
 \item Encontre os valores de $x$ no intervalo $[0,2\pi]$ para os quais $2\sin(x)=-1$.
 \item Qual é a soma das duas menores soluções positivas de $\mbox{sen}\left(3x-\frac{\pi}{4}\right)=0$?
 \item Encontre os valores de $x\in\mathbb{R}$ para os quais $\mbox{sen}(2x)=\cos x$.
 \item $2 \sen^2(x)- \sen(x)=0$, no intervalo $[0,2\pi]$.
 \item $2\sen^2(x)+\cos(x)=1$, no intervalo $[0,2\pi]$.
 \item $\tan^2(\theta) -3 = 0,$ no intervalo $[-\frac{\pi}{2},\frac{\pi}{2}]$.
\end{enumerate}
\item Encontre
\begin{enumerate}
    \item $\arcsin(\frac{\sqrt{3}}{2})$
    \item $\tan(\arccos(\frac{\sqrt{5}}{4}))$
    \item $\arctan(-\sqrt{3})$
    \item $\tan(\arctan(5))$
\end{enumerate}
\item Calcule o valor da expressão $\arcsen(-1)+\arccos(\frac{1}{2})-\arctg(\tg(\frac{\pi}{4}))$
\item Determine o domínio das seguintes funções
\begin{enumerate}
    \item $f(x)=\arcsin(x-3)$
    \item $g(x)=\arccos(x^2-1)$
\end{enumerate}
\item Use a Relação fundamental trigonométrica para mostrar que $\cos(\arcsin(x))=\sqrt{1-x^2}$, para $x \in [-1,1]$.
\end{enumerate}
\section{Gabarito:}
\begin{enumerate}
    \item $\tan(\alpha)= \frac{1}{3}$
    \item Sim, pois $\cos(\theta) = \frac{2}{3}$ e $\cos\left(\frac{\pi}{3}\right) = \frac{1}{2}$. Como $\cos(\theta) > \cos\left(\frac{\pi}{3}\right)$, então $\theta < \frac{\pi}{3}$, garantindo uma boa visão
\item
    \begin{enumerate}
        \item $\frac{5\pi}{6}$, $\sen\left(\frac{5\pi}{6}\right) = \frac{1}{2}$, $\cos\left(\frac{5\pi}{6}\right) = -\frac{\sqrt{3}}{2}$, $\tg\left(\frac{5\pi}{6}\right) = -\frac{\sqrt{3}}{3}$, $\sec\left(\frac{5\pi}{6}\right) = -\frac{2\sqrt{3}}{3}$, $\text{cossec}\left(\frac{5\pi}{6}\right) = 2$, $\text{cotg}\left(\frac{5\pi}{6}\right) = -\sqrt{3}$

        \item $\frac{7\pi}{4}$, $\sen\left(\frac{7\pi}{4}\right) = -\frac{\sqrt{2}}{2}$, $\cos\left(\frac{7\pi}{4}\right) = \frac{\sqrt{2}}{2}$, $\tg\left(\frac{7\pi}{4}\right) = -1$, $\sec\left(\frac{7\pi}{4}\right) = \sqrt{2}$, $\text{cossec}\left(\frac{7\pi}{4}\right) = -\sqrt{2}$, $\text{cotg}\left(\frac{7\pi}{4}\right) = -1$

        \item $\frac{7\pi}{6}$, $\sen\left(\frac{7\pi}{6}\right) = -\frac{1}{2}$, $\cos\left(\frac{7\pi}{6}\right) = -\frac{\sqrt{3}}{2}$, $\tg\left(\frac{7\pi}{6}\right) = \frac{\sqrt{3}}{3}$, $\sec\left(\frac{7\pi}{6}\right) = -\frac{2\sqrt{3}}{3}$, $\text{cossec}\left(\frac{7\pi}{6}\right) = -2$, $\text{cotg}\left(\frac{7\pi}{6}\right) = \sqrt{3}$
    \end{enumerate}
    \item $cos(x)=-\frac{\sqrt{7}}{4}, \tan(x)=\frac{3\sqrt{7}}{7}$ e $\sec(x)=-\frac{4\sqrt{7}}{7}$
    \item $\cos(\theta)= -\frac{\sqrt{15}}{4}$, $\tan(\theta)= -\frac{\sqrt{15}}{15}$ e $\sec(\theta)=-\frac{4\sqrt{15}}{15}$
    \item $\pm\frac{4}{5} $. Primeiro ou quarto quadrante
\item
    \begin{enumerate}
        \item $ $
            \begin{figure}[H]
                \centering
                \includegraphics[width=0.6\textwidth]{gabarito_aula5-7a}
                \caption{Gráfico da função $f(x) = \tan(x+\frac{\pi}{4})$ }
            \end{figure}

        \item $ $
            \begin{figure}[H]
                \centering
                \includegraphics[width=0.6\textwidth]{gabarito_aula5-7b}
                \caption{Gráfico da função $f(x) = \cos(4x)$ }
            \end{figure}
        \newpage
        \item $ $
            \begin{figure}[H]
                \centering
                \includegraphics[width=0.6\textwidth]{gabarito_aula5-7c}
                \caption{Gráfico da função $f(x) = 1-2\sin(x)$ }
            \end{figure}
        \item $ $
            \begin{figure}[H]
                \centering
                \includegraphics[width=0.6\textwidth]{gabarito_aula5-7d}
                \caption{Gráfico da função $f(x) = \sin(x)+2$ }
            \end{figure}
    \end{enumerate}
    \item
    \begin{enumerate}
        \item $-\cos(x)$
        \item $4\cos(x)$
        \item $\text{cossec}^2(x)$
    \end{enumerate}
    \item
\begin{enumerate}
    \item $-\text{cossec}(45^{\circ}) = -\sqrt{2}$

    \item $-\text{cotg}(30^{\circ}) = -\sqrt{3}$

    \item $-\tg\left(\frac{\pi}{4}\right) - (-\text{cotg}\left(\frac{\pi}{4}\right)) = -1 - (-1) = 0$

    \item $-\cos\left(\frac{\pi}{3}\right) = -\frac{1}{2}$

    \item $-\sin(90^{\circ}) = -1$

    \item $\cos\left(\frac{\pi}{6}\right) - \sen(0) = \frac{\sqrt{3}}{2} - 0 = \frac{\sqrt{3}}{2}$

    \item $-\tan(45^{\circ}) - \cos(30^{\circ}) = -1 - \frac{\sqrt{3}}{2} = \frac{-2-\sqrt{3}}{2}$
\end{enumerate}
    \item  Maior: $10$. Menor: $0$
    \item
    \begin{enumerate}
        \item Altura máxima $151m$ e altura mínima $1m$
        \item $30$min
    \newpage
        \item $ $
            \begin{figure}[H]
                \centering
                \includegraphics[width=0.6\textwidth]{gabarito_aula5-11c}
                \caption{Gráfico da função $h(t) = 76 + 75 \sin \left( \frac{\pi t}{15} - \frac{\pi}{2} \right)$ }
            \end{figure}
    \end{enumerate}
    \item $\sin(75^{\circ})= \frac{\sqrt{6}+ \sqrt{2}}{4}, \space \cos(75^{\circ})=\frac{\sqrt{6}-\sqrt{2}}{4}$ e $\tan(75^{\circ})= 2+\sqrt{3}$

\item
    \begin{enumerate}
        \item $\cos(3x) = 4\cos^3(x) - 3\cos(x)$
        \item $\cos(2x) = \cos^2(x) - \sen^2(x) = (1 - \sen^2(x)) - \sen^2(x) = 1 - 2\sen^2(x)$
        \item $\tg(2x) = \frac{\sen(2x)}{\cos(2x)} = \frac{2\sen(x)\cos(x)}{\cos^2(x) - \sen^2(x)} = \frac{2\tg(x)}{1 - \tg^2(x)}$
        \item $\sen(x-y) = -\frac{2\sqrt{5}}{25}, \cos(x-y) = \frac{11\sqrt{5}}{25}$ e $\tg(x-y) = -\frac{2}{11}$
        \item $\sen(15^{\circ}) = \frac{\sqrt{6}-\sqrt{2}}{4}$ e $\sen(7,5^{\circ}) = \frac{\sqrt{8-2\sqrt{6}-2\sqrt{2}}}{4}$
        \item $\cos(x)\cos\left(\frac{\pi}{2}\right) - \sen(x)\sen\left(\frac{\pi}{2}\right) = 0 - \sen(x) = -\sen(x)$
        \item $\ln(\sec\theta \cdot \cos\theta) = \ln(1) = 0$
        \item $\frac{\sqrt{2}-x}{\sqrt{3(x^2+1)}}$
    \end{enumerate}
\item $\sin(y)=\frac{7\sqrt{2}}{26}$ e $\cos(y)=\frac{17\sqrt{2}}{26}$
\item $-\tan(x)$
\item $\sin(2x)=2mn$, $\cos(2x)=n^2-m^2$ e $\tan(2x)=\frac{2mn}{n^2-m^2}$
\item $\sin(105^{\circ}) = \frac{\sqrt{6}+\sqrt{2}}{4}$
$\sec(105^{\circ}) = -\sqrt{2}-\sqrt{6}$
$\cos(15^{\circ}) = \frac{\sqrt{6}+\sqrt{2}}{4}$
$\tan(15^{\circ}) = 2-\sqrt{3}$
\item $\sin\left( \frac{x}{2}\right)=\sqrt{\frac{1}{2}\left(1+\frac{2\sqrt{2}}{3} \right)}, \,\, \cos\left( \frac{x}{2}\right)=\sqrt{\frac{1}{2}\left(1-\frac{2\sqrt{2}}{3} \right)}$ e $\tan\left( \frac{x}{2}\right)=3+2\sqrt{2}$
\item $ $
    \begin{figure}[H]
    \centering
    \includegraphics[width=0.6\textwidth]{gabarito_aula5-19}
    \caption{Gráfico de $g(x) = -1 + 2\cos(x + \frac{\pi}{2})$}
    \end{figure}
\item
\begin{itemize}
\item[a)] $x=\frac{7\pi}{6}$ ou $x=\frac{11\pi}{6}$
\item[b)]$\frac{\pi}{2}$
\item[c)]$x=\frac{\pi}{6}+2n\pi$, $x=\frac{5\pi}{6}+2n\pi$ ou $x=\frac{\pi}{2}+n\pi$,  onde $n\in\mathbb{Z}$
\item[d)]$x=0$, $x=\pi$, $x=2\pi$, $x = \frac{\pi}{6}$ ou  $x = \frac{5\pi}{6}$
\item[e)]$x=0$, $x= 2\pi$, $x=\frac{2\pi}{3}$ ou $x=\frac{4\pi}{3}$
\item[f)] $x = \pm\frac{\pi}{6}$
\end{itemize}

\item
\begin{enumerate}
    \item $\frac{\pi}{3}$
    \item $\frac{\sqrt{55}}{5}$
    \item $-\frac{\pi}{3}$
    \item $5$
\end{enumerate}
\item $-\frac{5\pi}{12}$
\item
    \begin{enumerate}
        \item $D_f= \{ x \in \mathbb{R} | 2 \leq x \leq 4\}$
        \item $D_g= \{ x \in \mathbb{R}| -\sqrt{2} \leq x \leq \sqrt{2}\}$
    \end{enumerate}

\item Seja $\theta = \arcsin(x) \Rightarrow \sin(\theta) = x$. Pela relação $\sin^2\theta + \cos^2\theta = 1$, temos $x^2 + \cos^2\theta = 1 \implies \cos\theta = \sqrt{1-x^2}$. Logo, $\cos(\arcsin(x)) = \sqrt{1-x^2}$
\end{enumerate}
